\begin{table}[htbp]
\centering
\caption{Salary and Rank Boosts for Non-Drafted Transferees into Draft-Vacancy Destinations}
\label{tab:misallocation-destination-draftee-wage}
\footnotesize
\begin{threeparttable}
\begin{tabular}{lccc}
\toprule
 & High-wage draftee & Low-wage draftee & Low $-$ High \\
\midrule
Transferees & 155 & 45 & -110 \\
With baseline salary & 20 & 9 & -11 \\
Current salary & 91.73 & 57.96 & -33.77 \\
Baseline salary & 84.40 & 90.56 & 6.16 \\
Salary change & -23.45 & -50.44 & -26.99 \\
Median salary change & 4.50 & -42.00 & -46.50 \\
Log salary change & -0.812 & -1.605 & -0.793 \\
Share with salary increase & 55.0\% & 22.2\% & -32.8\% \\
\midrule
Position-rank change & 0.071 & 0.156 & 0.085 \\
Share with position-rank boost & 7.7\% & 26.7\% & 18.9\% \\
Court-rank change & -0.123 & -0.178 & -0.055 \\
Share with court-rank boost & 0.6\% & 0.0\% & -0.6\% \\
\bottomrule
\end{tabular}
\begin{tablenotes}[flushleft]\footnotesize
\item \textit{Notes:} Sample consists of non-drafted workers who transferred to a draft-vacancy destination office-year with observed drafted-worker salary, using the same destination wage split as Model G. High-wage draftee destinations are office-years where the mean salary of drafted workers is at or above the median among draft-vacancy office-years with observed drafted-worker salary; low-wage destinations are below that median. Salary changes compare current-year salary to the worker's own baseline (lagged) salary, so those rows use only transferees with observed baseline salary. Position rank is coded 1--3 from title hierarchy; court rank is coded 0--8 with higher values indicating higher status. Rank-change rows code missing ranks as 0, matching the MNL convention.
\end{tablenotes}\end{threeparttable}\end{table}
